% Minimal retained preamble for lane-p2 cards taken from
% sources/logic-fl02-conservation-of-ramsey-s-theorem-f/pinned-original-source.tex
% (Conservation of Ramsey's theorem for pairs, arXiv:2402.11616v3).
%
% Copied verbatim from the author's own file:
%   line 2      : \documentclass{article}
%   lines 16-25 : the author's theorem declarations
%   lines 46-88 : the author's macro definitions for the reverse-mathematics principles
%   line 91     : \card
%   lines 104-107: the natural sums and products, which rely on mathabx (source line 103)
%   line 116    : \qt
%
% Omitted author preamble lines and why (nothing was renamed or redefined):
%   line 3 (graphicx), 7 (url), 9 (cancel): no figure, url or cancel use in the retained spans;
%   line 6 (cleveref): the assembler adds its own cleveref machinery, a second load clashes;
%   line 8 (\usepackage[utf8x]{inputenc}): inputenc with the utf8x option fights the forced
%       xelatex engine;
%   lines 28-37 (the rep@theorem machinery and \newreptheorem{maintheorem}): the retained spans
%       use no rep-theorem environment;
%   lines 40-43 (xcolor and the three author comment macros): not used by the retained spans;
%   lines 93-99 (the LS2 font encoding block for \smashtimes based on stix): the host TeX Live has
%       no stix directory and the author himself never uses the result -- line 105 maps \ntimes to
%       \dottimes, and line 110 shows \smashtimes already commented out in the source.
% mathabx is retained because \dotplus, \dottimes and \dotdiv come from it.
\documentclass{article}
\usepackage{amsmath,amssymb,amsthm}
\usepackage{mathabx}
\usepackage{cancel}
\usepackage{cleveref}

\newtheorem{theorem}{Theorem}
\numberwithin{theorem}{section}
\newtheorem{maintheorem}[theorem]{Main Theorem}
\newtheorem{lemma}[theorem]{Lemma}
\newtheorem{question}[theorem]{Question}
\newtheorem{proposition}[theorem]{Proposition}
\newtheorem{statement}[theorem]{Statement}
\newtheorem{remark}[theorem]{Remark}
\newtheorem{definition}[theorem]{Definition}
\newtheorem{corollary}[theorem]{Corollary}

\newcommand{\RCA}{\mathsf{RCA}}
\newcommand{\WKL}{\mathsf{WKL}}
\newcommand{\SRT}{\mathsf{SRT}}
\newcommand{\EM}{\mathsf{EM}}
\newcommand{\SEM}{\mathsf{SEM}}
\newcommand{\ADS}{\mathsf{ADS}}
\newcommand{\SCAC}{\mathsf{SCAC}}
\newcommand{\SADS}{\mathsf{SADS}}
\newcommand{\COH}{\mathsf{COH}}
\newcommand{\ACA}{\mathsf{ACA}}
\newcommand{\qvdash}{\operatorname{{?}{\vdash}}}
\newcommand{\nqvdash}{\operatorname{{?}{\nvdash}}}
\newcommand{\BSig}{\mathsf{B}\Sigma^0}
\newcommand{\ISig}{\mathsf{I}\Sigma^0}
\newcommand{\LSig}{\mathsf{L}\Sigma^0}
\newcommand{\BPi}{\mathsf{B}\Pi^0}
\newcommand{\IPi}{\mathsf{I}\Pi^0}
\newcommand{\LPi}{\mathsf{L}\Pi^0}
\newcommand{\IDel}{\mathsf{I}\Delta^0}
\newcommand{\BCDe}{\mathsf{BC}\Delta^0}
\newcommand{\WF}{\mathsf{WF}}
\newcommand{\RT}{\mathsf{RT}}
\newcommand{\BME}{\mathsf{BME}}
\newcommand{\TT}{\mathsf{TT}}
\renewcommand{\O}{\mathcal{O}}
\renewcommand{\L}{\mathcal{L}}
\newcommand{\C}{\mathcal{C}}
\newcommand{\R}{\mathcal{R}}
\renewcommand{\S}{\mathcal{S}}
\newcommand{\T}{\mathcal{T}}
\renewcommand{\P}{\mathcal{P}}
\newcommand{\M}{\mathcal{M}}
\newcommand{\Nc}{\mathcal{N}}
\renewcommand{\P}{\mathcal{P}}
\newcommand{\Q}{\mathcal{Q}}
\newcommand{\NN}{\mathbb{N}}
\newcommand{\PP}{\mathbb{P}}
\newcommand{\uh}[0]{{\upharpoonright}}
\newcommand{\card}{\operatorname{card}}
\newcommand{\nplus}{\dotplus}
\newcommand{\ntimes}{\dottimes}
\newcommand{\nminus}{\dotdiv}
\newcommand{\nsum}{\dot{\sum}}
\def\qt#1{``#1''}
\title{A valid node of the SEM forcing yields an extension of the underlying condition}
\author{Quentin Le Houérou, Ludovic Levy Patey, Keita Yokoyama}
\date{arxiv:2402.11616v3; CC BY 4.0}
\begin{document}
\maketitle
\noindent\textit{Editorial note.} The counted unit is the article's own Lemma lem:sem-wf-valid-leaf-force-positive at source lines 841-844 together with the proof environment that begins on the next line (845-850); the pairing is the ordinary adjacent one.\par\medskip
\noindent\textit{Editorial note.} The retained context at source lines 819-829 is the author's own definition of a node of a condition, of the union pi of a node, of validity and of satisfaction of a Sigma\textasciicircum{}0\_1 formula; the retained context at source lines 831-833 is the statement of the immediately preceding transitivity lemma, which the counted proof invokes and which is counted by a sibling card of this lane.\par\medskip
\noindent\textit{Editorial note.} The retained preamble reproduces verbatim the author's document class at source line 2, his theorem declarations at source lines 16-25, his macro definitions at source lines 46-88 and 91, his natural-operation macros at source lines 104-107 and his quotation macro at source line 116, and it retains mathabx from source line 103 and cleveref, which the assembler only adds by itself when it generates navigation stubs. The LS2 font-encoding block at source lines 93-99 is not retained because the host TeX Live has no stix directory, and the inputenc load with the utf8x option at source line 8 is not retained because it fights the forced xelatex engine.\par\medskip
\noindent\textit{Editorial note.} No retained line contains a citation, so the material carries no bibliography; the only cross-reference of the retained proof is the author Cref call to the transitivity lemma, whose statement is retained in the context.\par\medskip
\noindent\textit{Editorial note.} The article's conservation theorems, its COH and WKL lemmas and the other lemmas of the SEM section are not counted here and are not used as premises.\par\medskip
\noindent\textit{Allowed standard prerequisites (editorial).}\begin{itemize}
\item Node of a condition, validity and the set pi of a node. A (sigma, a)-node is an M-finite sequence of pairs (i\_s, rho\_s) with each rho\_s a non-empty M-finite R-transitive subset of X\_a, increasing in the order of the integers and comparable in the direction prescribed by the bit i\_s; it is valid when every block is contained in the corresponding side A\_\{i\_s\}(R) of the colouring; pi of the node is the union of sigma with all blocks; the author's own definition is retained at source lines 819-829.. Fixes the objects of the counted lemma, namely a valid node and its union pi, whose transitivity is what makes the pair (pi, b) a condition.
\item Stability of the colouring R and the collection scheme. Stability of R means that for every x in X\_a the set of y above x on which R holds is either finite or cofinite, which turns the validity of a node into infinitary comparability statements; the collection scheme for Sigma\textasciicircum{}0\_2 formulas then lets the finitely many bounds obtained that way be taken uniformly.. The proof converts validity into the statements that each block is below or above almost all of X\_a, and then uses collection to find a single initial segment whose removal makes the comparisons uniform.
\item The set attached to a node is transitive. For every node the union of sigma with all its blocks is R-transitive; this is the statement of the immediately preceding lemma, retained with this card at source lines 831-833.. The proof invokes it to conclude that the set tau it constructs is R-transitive, which is one of the two conditions for (tau, b) to be an extension of (sigma, a).
\end{itemize}\medskip
\begin{definition}
Fix a condition $(\sigma, a)$.
\begin{enumerate}
    \item[(1)] A \emph{$(\sigma, a)$-node} is an $M$-finite sequence of the form
$$\nu = \langle (i_0, \rho_0), (i_1, \rho_1), \dots (i_{q-1}, \rho_{q-1}) \rangle$$ such that for every~$s < q$, $i_s < 2$, $\rho_s \subseteq X_a$ is a non-empty $M$-finite $R$-transitive set with $\min \rho_s > \max \rho_{s-1}$ (letting $\rho_{-1} = \sigma$), and for every~$s < t < q$, if $i_s = 0$ then $\rho_s \to_R \rho_t$, and if $i_s = 1$ then $\rho_t \to_R \rho_s$.
We then write $$\pi(\nu) = \sigma \cup \rho_0 \cup \dots \cup \rho_{q-1}$$
    \item[(2)] This $(\sigma, a)$-node is \emph{valid} if for every $s < q$, $\rho_s \subseteq A_{i_s}(R)$
%    \item[(3)] A (finite or infinite) set~$Y$ is \emph{compatible} with this $(\sigma, a)$-node if for every if for every~$s < q$, if $i_s = 0$ then $\rho_s \to_R Y$, and if $i_s = 1$ then $Y \to_R \rho_s$.
    \item[(3)] A $(\sigma, a)$-node $\nu$ \emph{satisfies} a $\Sigma^0_1(\M)$-formula $\varphi(G, x) \equiv \exists i \psi(G \uh i, x)$  if $\psi(\pi(\nu), |\nu|-1)$ holds.
\end{enumerate}
\end{definition}

\begin{lemma}\label[lemma]{lem:sem-wf-nodes-transitive}
Fix a $(\sigma, a)$-node $\nu$. Then $\pi(\nu)$ is $R$-transitive.
\end{lemma}

\begin{lemma}\label[lemma]{lem:sem-wf-valid-leaf-force-positive}
Fix a valid $(\sigma, a)$-node~$\nu$.
Then $(\pi(\nu), b)$ is an extension of~$(\sigma, a)$ for some~$b \in M$.
\end{lemma}

\begin{proof}
Say $\nu = \langle (i_0, \rho_0), (i_1, \rho_1), \dots (i_{q-1}, \rho_{q-1}) \rangle$.
By definition of validity and stability of~$R$, for every~$s < q$, if $i_s = 0$ then $\forall x \in \rho_s \forall^\infty y \in X_a$, $R(x, y)$, and if $i_s = 1$ then $\forall x \in \rho_s \forall^\infty y \in X_a$, $R(y, x)$.
By $\BSig_2$, there exists some~$r \in M$ such that for every~$s < q$, if $i_s = 0$ then $\rho_s \to_R X_a \setminus \{0, \dots, r\}$, and if $i_\ell = 1$ then $X_a \setminus \{0, \dots, r\} \to_R \rho_s$. Let~$b \in M$ be such that $X_b = X_a \setminus \{0, \dots, r\}$ and let~$\tau = \sigma \cup \rho_0 \cup \dots \cup \rho_{t-1}$.
By \Cref{lem:sem-wf-nodes-transitive}, $\tau$ is $R$-transitive and by construction, $X_b \subseteq X_a$ is included in a minimal interval of $\tau$. Thus $(\tau, b)$ is an extension of $(\sigma, a)$.
\end{proof}

\end{document}
